El presente trabajo tuvo como alcance la reingeniería mecánica parcial y la automatización de un brazo robótico cartesiano para su integración en una línea de empaquetado del Laboratorio de Diseño de Procesos de la Universidad del Valle de Guatemala. El proyecto se enfocó en mejorar el desempeño mecánico del sistema mediante el rediseño de los mecanismos de traslación de los ejes X, Y y Z, la incorporación de ejes guía de acero y la integración de rodamientos lineales.

Asimismo, se contempló la automatización del brazo para ejecutar tareas de manipulación y transferencia de productos mediante una secuencia de tipo \textit{pick and place}. El sistema localizó las piezas sobre la banda transportadora mediante visión computacional, a partir de su forma, con un modelo de segmentación entrenado para tipos de pieza definidos. Con base en esa posición y en el avance de la banda medido por un encoder, el brazo recogió cada pieza en una estación de captura fija y la trasladó a un punto de entrega.

El proyecto también incluyó un modo de operación manual mediante un control inalámbrico, que permitió mover los ejes del brazo de forma independiente durante tareas de posicionamiento, pruebas, mantenimiento o intervención del operador. Además, se integró una rutina de referenciación automática apoyada en finales de carrera para establecer las referencias de movimiento, y una pantalla OLED como interfaz local para visualizar el modo de operación, el estado de las calibraciones, el estado de la cámara y las fases del ciclo automático.

Sin embargo, el proyecto se limitó a la automatización del brazo robótico cartesiano y no contempló el diseño completo de una línea de empaquetado industrial. Las pruebas se realizaron en un entorno académico y controlado, por lo que los resultados estuvieron condicionados por las características del laboratorio, los recursos disponibles y las dimensiones del prototipo. En consecuencia, el sistema no fue validado bajo condiciones industriales de producción continua, cargas elevadas o ambientes variables.

El sistema de visión se limitó a las piezas para las que se entrenó el modelo; el firmware aceptó únicamente los tipos de pieza 6 y 7. No se desarrolló un sistema destinado a identificar cualquier tipo de objeto, material o geometría, y su desempeño estuvo condicionado por la iluminación, la posición de la cámara, el contraste entre las piezas y el fondo y la estabilidad del entorno visual. El ciclo automático atendió una pieza a la vez: las escenas con varias piezas se descartaron. Todas las piezas se trasladaron al mismo punto de entrega, por lo que el tipo detectado no se utilizó para asignar destinos distintos. Además, la cámara entregó únicamente la caja delimitadora de cada detección, por lo que el sentido de giro de las piezas diagonales no pudo determinarse.

Finalmente, aunque el rediseño mecánico estuvo orientado a mejorar la fluidez del movimiento, el brazo permaneció condicionado por la estructura base existente, la capacidad de los actuadores, el espacio de trabajo disponible y las tolerancias de fabricación. La posición de los ejes se estimó mediante el conteo de los pasos ordenados, sin una medición independiente de la posición del efector final. La validación se centró en la tasa de recogida y entrega del ciclo automático; la exactitud y la repetibilidad de posicionamiento del efector final quedaron fuera de las mediciones realizadas.
